%% ============================================================================
%%  <과목 이름> --- 문제풀이, 국문판.
%%  빌드: ./build.sh problems   (latexmk -xelatex, 이어서 tools/probcheck.py)
%%
%%  problems-en.tex 와 장·문제의 차례와 번호가 같아야 한다.
%% ============================================================================
\documentclass[11pt]{article}
\usepackage[korean]{coursenotes}

\coursename{샘플 과목}
\term{2027년 2학기}
\lectitle{문제와 풀이}
\lecdate{2027년 9월 6일 기준}

\newcommand{\ProbWord}{문제}
\newcommand{\SolWord}{풀이}
\newcommand{\SourceWord}{출처}
\newcommand{\NotesWord}{노트}

%% ============================================================================
%%  problems/common.tex --- shared by problems-en.tex and problems-ko.tex
%%
%%  The problem book alternates problem and solution: every problem starts a
%%  new page, and its solution starts the next new page.  The web viewer
%%  (pdf-notes-viewer, DOCS entry with solutions:true) covers the solution
%%  pages until they are tapped.  It finds them through named destinations
%%  that \problem and \solution put into the PDF:
%%      pb.prob.<n.m>   where problem n.m starts
%%      pb.sol.<n.m>    where its solution starts
%%      pb.end          where back matter after the last solution starts
%%  A solution is covered from its pb.sol page up to the page before the next
%%  mark, or to the last page.  Anything placed after the last solution (an
%%  appendix, say) must therefore start on a new page with \pbmark{end}.
%%  "pb" must equal MARK in the viewer's configuration block.
%%
%%  The labels chapstart:<n>, probstart:<n.m>, solstart:<n.m> are for
%%  tools/probcheck.py, which checks that the two editions match.
%%
%%  Numbering: chapters use the section counter, problems the subsection
%%  counter, so bookmarks read "2 Steady states" and "2.1 A membrane (...)":
%%  the viewer keys on that leading number to keep your place across the
%%  EN/KO switch.  Both editions must carry the same chapters and problems in
%%  the same order.  Equations are numbered afresh in every problem.
%%
%%  Each main file defines the words used below:
%%    \ProbWord \SolWord \SourceWord \NotesWord
%%  Requires the course style (coursenotes.sty: hyperref, enumitem, fancyhdr).
%% ============================================================================

%% ---- notation shared by the problems (add the course's own here) ----------
\newcommand{\ans}[1]{\boxed{\;#1\;}}
\newcommand{\e}{\mathrm{e}}

%% ---- problem parts ---------------------------------------------------------
\newenvironment{parts}%
  {\begin{enumerate}[label=(\alph*),leftmargin=2.2em,itemsep=4pt,topsep=4pt]}%
  {\end{enumerate}}
%% a part of a solution: \solpart{a}{Short title}
\newcommand{\solpart}[2]{%
  \par\addvspace{9pt}\noindent\textbf{(#1)\enspace#2}\par\nobreak\vspace{2pt}}

%% ---- equations restart in every problem ------------------------------------
\makeatletter
\@addtoreset{equation}{subsection}
\renewcommand{\theHequation}{\thesubsection.\arabic{equation}}
\makeatother

%% ---- chapters, problems, solutions ----------------------------------------
\makeatletter
\newif\ifpb@afterchap
\def\pb@probtitle{}

%% \pbmark{name}: a named destination pb.<name> on the current page, for the
%% web viewer (see the top of this file).  Same mechanism as \phantomsection.
\newcommand{\pbmark}[1]{\Hy@raisedlink{\hyper@anchorstart{pb.#1}\hyper@anchorend}}

%% \probchapter{Title}{notes range}
\newcommand{\probchapter}[2]{%
  \clearpage
  \refstepcounter{section}%
  \phantomsection
  \addcontentsline{toc}{section}{\protect\numberline{\thesection}#1}%
  \edef\pb@id{\thesection}%
  \label{chapstart:\pb@id}%
  \markright{#1}%
  \vspace*{-0.2\baselineskip}%
  \noindent\rule{\textwidth}{1.1pt}\par\nobreak\vspace{7pt}%
  \noindent{\Large\bfseries\thesection\quad #1}\par\nobreak\vspace{3pt}%
  \noindent{\small\itshape #2}\par\nobreak\vspace{5pt}%
  \noindent\rule{\textwidth}{0.4pt}\par\nobreak
  \vspace{16pt}%
  \global\pb@afterchaptrue}

%% \problem{Title}{source}{related notes}
\newcommand{\problem}[3]{%
  \ifpb@afterchap\global\pb@afterchapfalse\else\clearpage\fi
  \refstepcounter{subsection}%
  \gdef\pb@probtitle{#1}%
  \phantomsection
  \addcontentsline{toc}{subsection}%
    {\protect\numberline{\thesubsection}#1\texorpdfstring{\ \ {\normalfont\small(#2)}}{ (#2)}}%
  \edef\pb@id{\thesubsection}%
  \label{probstart:\pb@id}%
  \pbmark{prob.\pb@id}%
  \markright{\ProbWord\ \thesubsection}%
  \noindent{\large\bfseries\ProbWord\ \thesubsection\quad #1}\par\nobreak
  \vspace{3pt}%
  \noindent{\small\SourceWord: #2\hfill\NotesWord: #3}\par\nobreak
  \vspace{2pt}\noindent\rule{\textwidth}{0.4pt}\par\nobreak
  \vspace{6pt}}

%% \solution   (starts the solution of the current problem on a new page)
\newcommand{\solution}{%
  \clearpage
  \phantomsection
  \edef\pb@id{\thesubsection}%
  \label{solstart:\pb@id}%
  \pbmark{sol.\pb@id}%
  \markright{\SolWord\ \thesubsection}%
  \noindent{\large\bfseries\SolWord\ \thesubsection\quad\pb@probtitle}\par\nobreak
  \vspace{2pt}\noindent\rule{\textwidth}{0.4pt}\par\nobreak
  \vspace{6pt}}

%% the running head: title on the left, the current problem/solution on the right
\newcommand{\probheads}{%
  \fancyhead[R]{\small\rightmark}}
\makeatother

%% ---- table of contents: problems are listed under their chapter -------------
\setcounter{tocdepth}{2}


\begin{document}
\makelecheader
\probheads

\tableofcontents

\clearpage
\markboth{}{}
\section*{이 문제집을 쓰는 법}
문제는 늘 새 쪽에서 시작하고 그 풀이는 다음 쪽에서 시작한다. 쪽을 넘기기 전에
먼저 풀어 볼 수 있다. 웹 뷰어에서는 풀이 쪽이 가려져 있다가 누르면 드러난다.

\emph{출처.} 제목 아래 출처 줄에 문제가 어디서 왔는지 적었다. 예년 시험이나
과제(연도와 함께), 이 노트에 맞춰 만든 모의고사, 이 문제집을 위해 새로 만든
문제(``새 문제'')다. 예년 문제에 문항을 덧붙였으면 ``+ 새 문항''이라고 적었다.

\emph{기호}는 강의노트를 따른다. ``노트'' 칸은 그 문제가 기대는 절이다.

%% ============================================================================
\probchapter{무작위 걷기와 확산}{노트 \S1--\S2}

%% ----------------------------------------------------------------------------
\problem{치우친 걷기}{새 문제}{\S1}

직선 위의 걷는 이가 서로 독립인 걸음 $s_k$ 를 내딛는다. 걸음은 확률 $p$ 로
$+a$, 확률 $q=1-p$ 로 $-a$ 다. $n$ 걸음 뒤의 위치는 $x_n=\sum_{k=1}^n s_k$ 다.
\begin{parts}
  \item 한 걸음의 평균과 분산을 구하라.
  \item $\avg{x_n}$ 과 $x_n$ 의 분산을 구하라.
  \item 표류 $\avg{x_n}$ 은 $n$ 에, 퍼짐은 $\sqrt n$ 에 비례해 자란다.
        $p=0.6$ 일 때 몇 걸음부터 표류가 표준편차보다 커지는가?
\end{parts}

\solution

\solpart{a}{한 걸음}
\begin{equation}
  \avg{s}=a\,p-a\,q=a(2p-1),\qquad \avg{s^2}=a^2
\end{equation}
이므로
\begin{equation}
  \operatorname{Var}s=\avg{s^2}-\avg{s}^2=a^2\bigl[1-(2p-1)^2\bigr]
  =a^2\bigl[(1-2p+1)(1+2p-1)\bigr]=\ans{4a^2pq}.
\end{equation}
확인: $p=\tfrac12$ 이면 $a^2$ 로 치우치지 않은 걷기, $p=1$ 이면 $0$ 으로 정해진
걷기가 된다.

\solpart{b}{$n$ 걸음 뒤}
평균은 언제나 더해지고, 분산은 걸음이 독립이므로 더해진다.
\begin{equation}
  \ans{\avg{x_n}=na(2p-1),\qquad \operatorname{Var}x_n=4na^2pq}.
\end{equation}

\solpart{c}{표류와 퍼짐}
표류가 표준편차보다 커지는 조건은
\begin{equation}
  na\,\lvert 2p-1\rvert>2a\sqrt{npq}
  \iff n>\frac{4pq}{(2p-1)^2}
\end{equation}
이다. $p=0.6$ 이면 $4pq=0.96$, $(2p-1)^2=0.04$ 이므로 $n>24$, 곧
\ans{\text{25번째 걸음부터}} 다.

%% ----------------------------------------------------------------------------
\problem{확산에는 얼마나 걸리나}{새 문제}{\S2}

물속의 작은 분자는 $D=1.0\times10^{-9}\,\mathrm{m^2\,s^{-1}}$ 이다.
일차원 결과 $\avg{x^2}=2Dt$ 를 쓴다.
\begin{parts}
  \item 제곱평균제곱근 변위가 $1\,\mathrm{mm}$ 가 되기까지 얼마나 걸리나?
  \item 세포 하나($10\,\mu\mathrm{m}$)를 건너는 데는 얼마나 걸리나?
  \item 두 시간의 비를 한 줄로 설명하라.
\end{parts}

\solution

\solpart{a}{1 밀리미터}
\begin{equation}
  t=\frac{\avg{x^2}}{2D}=\frac{(10^{-3}\,\mathrm m)^2}{2\times10^{-9}\,\mathrm{m^2\,s^{-1}}}
  =\ans{500\,\mathrm s}\approx 8\,\text{분}.
\end{equation}

\solpart{b}{10 마이크로미터}
\begin{equation}
  t=\frac{(10^{-5}\,\mathrm m)^2}{2\times10^{-9}\,\mathrm{m^2\,s^{-1}}}
  =\ans{0.05\,\mathrm s}.
\end{equation}

\solpart{c}{비}
시간은 거리의 제곱에 비례한다. 거리가 $100$ 배면 시간은 $100^2=10^4$ 배이고,
실제로 $500\,\mathrm s/0.05\,\mathrm s=10^4$ 이다.

%% ============================================================================
\probchapter{정상 상태}{노트 \S3}

%% ----------------------------------------------------------------------------
\problem{막을 지나는 흐름}{새 문제}{\S3}

두께 $L$ 인 막이 잘 섞인 두 용액을 가른다. 막 안의 농도는 확산 방정식을
따르고, $c(0,t)=c_0$ 와 $c(L,t)=c_L$ 은 고정되어 있다.
\begin{parts}
  \item 정상 상태의 분포 $c(x)$ 를 구하라.
  \item 정상 상태의 흐름 $J=-D\,\partial c/\partial x$ 를 구하라.
  \item $D=1.0\times10^{-9}\,\mathrm{m^2\,s^{-1}}$,
        $L=100\,\mu\mathrm m$, $c_0=1\,\mathrm{mol\,m^{-3}}$, $c_L=0$ 일 때
        $J$ 의 값을 구하라.
\end{parts}

\solution

\solpart{a}{분포}
정상 상태에서는 $\partial c/\partial t=0$ 이므로 확산 방정식에서
$\d^2c/\d x^2=0$ 만 남는다. 분포는 직선이다. 두 경계값에 맞추면
\begin{equation}
  \ans{c(x)=c_0+(c_L-c_0)\,\frac xL}.
\end{equation}

\solpart{b}{흐름}
\begin{equation}
  J=-D\,\dv{c}{x}=\ans{D\,\frac{c_0-c_L}{L}}
\end{equation}
이고 어느 $x$ 에서나 같다. 정상 상태라면 어디에도 물질이 쌓이지 않아야 하므로
당연하다. 흐름은 농도가 높은 쪽에서 낮은 쪽으로 간다.

\solpart{c}{값}
\begin{equation}
  J=\frac{(1.0\times10^{-9}\,\mathrm{m^2\,s^{-1}})(1\,\mathrm{mol\,m^{-3}})}{1.0\times10^{-4}\,\mathrm m}
  =\ans{1.0\times10^{-5}\,\mathrm{mol\,m^{-2}\,s^{-1}}}.
\end{equation}


\end{document}
